\section{Methodology}

\subsection{Corpora}

We use three corpora spanning acted and spontaneous speech from children and adults:

\begin{itemize}
    \item \textbf{C-BESD (MY)}: 4,179 utterances from 70 children (ages 6--12), six emotion classes, bilingual (English and Telugu). Acted/portrayed speech.
    \item \textbf{FAU Aibo}: 18,216 utterances from 51 children (ages 10--13), four emotion classes (angry, happy, neutral, sad). Spontaneous child-robot interaction.
    \item \textbf{IEMOCAP}: $\sim$9,794 utterances from 10 adults, four emotion classes (angry, happy, neutral, sad). Acted/portrayed speech, serving as adult control.
\end{itemize}

For cross-corpus compatibility, we map all corpora to four shared emotion classes (angry, happy, neutral, sad) except for C-BESD in-domain evaluations where six native classes are used.

\subsection{Speaker-Independent Split}

We use a deterministic MD5-hash-based speaker split (70/15/15 train/val/test) with a fixed data split seed of 42. This ensures zero speaker leakage between splits and reproducibility across all experiments.

\subsection{Model Architecture}

Our model consists of four components built on a frozen or unfrozen WavLM Base backbone:

\begin{enumerate}
    \item \textbf{SSL Backbone}: WavLM Base (microsoft/wavlm-base-sv, 94M parameters) extracts frame-level 768-dimensional features at 50 Hz from 16kHz audio.
    \item \textbf{Layer Fusion}: A 12-layer learnable weighted summation module combines hidden states from all transformer layers. When disabled, the last layer alone is used.
    \item \textbf{Temporal Pooling}: Three strategies are compared---mean pooling, self-attention pooling (a learned query vector attends over the time dimension), and prosody-guided pooling (F0 and energy features modulate the attention weights).
    \item \textbf{SEMLP Classifier}: A Squeeze-Excitation MLP head ($\sim$593K parameters) maps pooled features to class logits.
\end{enumerate}

Total trainable parameters are approximately 704K when the backbone is frozen. An optional Acoustic Calibration Adapter (a small bottleneck MLP inserted after layer fusion) is evaluated in the ablation study.

\subsection{Training Protocol}

All models are trained end-to-end on raw waveforms with a batch size of 16, AdamW optimizer, and 100-epoch maximum with early stopping (patience=15 on validation weighted accuracy). Each experiment is repeated with three random seeds (42, 123, 456) for reliability estimation.

\subsection{Evaluation Metrics}

We report weighted accuracy (WA) and unweighted average recall (UAR). Distribution shift is quantified using Fr\'{e}chet Distance (FD) and Subspace Mean Maximum Discrepancy (SMMD) between corpus pairs.

\subsection{Experimental Design}

Seven experimental phases (B1--B7, 192 total runs) systematically isolate each factor:

\begin{itemize}
    \item \textbf{B1 (Pooling Baseline)}: 3 pooling $\times$ 3 corpora $\times$ 3 seeds = 27 runs.
    \item \textbf{B2 (Zero-Shot Transfer)}: 18 cross-corpus directions, single run each.
    \item \textbf{B3 (Augmentation Sensitivity)}: 4 augmentation conditions $\times$ 3 corpora $\times$ 3 seeds = 36 runs.
    \item \textbf{B4 (Layer Fusion Ablation)}: 54 runs varying fusion mode and layer selection.
    \item \textbf{B5 (Unfreeze Comparison)}: 3 corpora $\times$ 3 seeds = 9 runs.
    \item \textbf{B6 (Module Ablation)}: 5 configurations $\times$ 2 corpora $\times$ 3 seeds = 30 runs.
    \item \textbf{B7 (Transfer Fine-tuning)}: 6 transfer directions $\times$ 3 seeds = 18 runs.
\end{itemize}
